\chapter*{Обзор литературы}
\addcontentsline{toc}{chapter}{Обзор литературы}

Настоящий обзор литературы выполняет рамочную функцию и не дублирует подробные обзоры, размещенные в последующих главах. В главе~\ref{ch:market} рассматриваются исследования музыкальной атмосферы, рынок музыки для бизнеса, существующие сервисные решения и правовая специфика публичного воспроизведения. В главах~\ref{ch:sequential}--\ref{ch:experiment} последовательно анализируются персональные и последовательные рекомендательные системы, методы групповых рекомендаций и экспериментальная проверка моделей на датасете YAMBDA. Поэтому общая логика работы строится от общего к частному: сначала рассматривается широкий контекст музыкальных рекомендаций и применения музыки в коммерческих пространствах, затем выделяется задача групповых музыкальных рекомендаций, после чего проверяется конкретная алгоритмическая гипотеза и описывается возможный сервисный сценарий.

Исследования музыкального сопровождения заведений показывают, что музыка является частью клиентского опыта, а не только нейтральным фоном. В работах по атмосферным эффектам и \textit{musicscape} подчеркивается связь музыкального оформления с эмоциональной реакцией посетителей, восприятием сервиса и поведенческими намерениями~\cite{roschk2017atmospheric,trompeta2022music}. При этом важны не просто наличие музыки или ее темп, а соответствие звучания контексту и субъективная привлекательность музыки для аудитории~\cite{caldwell2002tempo,demoulin2011music}. Этот вывод напрямую связан с постановкой настоящей работы: если эффект музыкального сопровождения зависит от того, насколько музыка подходит посетителям и ситуации, то возникает потребность в более гибкой модели выбора треков, чем заранее заданный статичный плейлист. Однако литература о музыкальной атмосфере в основном объясняет влияние музыки на опыт посетителя, но не описывает, как проектировать сервис, автоматически учитывающий предпочтения фактически присутствующей группы людей.

С технической стороны работа опирается на развитие рекомендательных систем. Классические методы коллаборативной фильтрации и матричной факторизации сформировали базовый подход к моделированию пользовательских предпочтений~\cite{koren2009matrix,rendle2009bpr}. Более современные последовательные модели, включая GRU4Rec, Caser, SASRec и BERT4Rec, учитывают порядок взаимодействий и краткосрочный контекст пользователя~\cite{hidasi2016session,tang2018caser,sasrec2018,bert4rec2019}. Для музыкального домена эта линия особенно важна, поскольку история прослушиваний содержит устойчивые паттерны повторного потребления, а крупные датасеты вроде YAMBDA позволяют проверять такие модели на масштабных данных~\cite{abbattista2024sequential,yambda2025}. В настоящей работе персональный последовательный скорер рассматривается не как конечная цель, а как нижний уровень системы: он формирует индивидуальные оценки треков, которые затем используются для построения коллективной рекомендации.

Следующий уровень анализа связан с групповыми рекомендательными системами. В литературе предложены как классические стратегии агрегации индивидуальных предпочтений, так и нейросетевые модели группового консенсуса~\cite{masthoff2004group,masthoff2011handbook,dara2020survey,felfernig2018handbook}. Модели AGREE и GroupIM являются важными точками сравнения, поскольку переходят от ручных правил агрегации к обучаемым представлениям группы~\cite{cao2018agree,sankar2020groupim}. Вместе с тем многие существующие подходы рассматривают группу как относительно устойчивый объект или исследуют задачу в абстрактной постановке, без учета музыкального домена, публичного пространства, правовой модели и бизнес-ограничений. Поэтому их можно использовать как алгоритмические бейзлайны, но нельзя напрямую считать готовыми решениями для заведений и мероприятий.

Практические B2B-сервисы музыки для бизнеса закрывают другую часть задачи: они предлагают легальное воспроизведение, централизованное управление и подбор плейлистов под тип заведения~\cite{zvuk_business_site,yandex_station_business}. Сервисы заказа музыки посетителями, напротив, дают гостям прямую интерактивность, но обычно передают влияние отдельному пользователю, а не формируют сбалансированную коллективную рекомендацию~\cite{gusli_site,bitzabid_site}. В результате между академическими исследованиями групповых рекомендаций и реальными сервисами музыкального сопровождения остается заметный разрыв. Именно практическому применению и проектированию сервисов групповых музыкальных рекомендаций для заведений и мероприятий в литературе уделено мало внимания: чаще исследуется либо влияние музыки на посетителей, либо алгоритмическая точность моделей, либо операционная сторона B2B-воспроизведения.

Настоящая работа позиционируется в этом промежутке. Она соединяет рыночную и правовую постановку задачи, концепцию сервиса и экспериментальную проверку алгоритмов коллективной рекомендации. В отличие от работ, ограниченных персональными рекомендациями, здесь рассматривается эфемерная группа посетителей. В отличие от большинства работ по групповым рекомендациям, внимание уделяется музыкальному домену, аудиоэмбеддингам, публичному пространству и ограничениям бизнеса. В отличие от существующих B2B-решений, предлагается не только управляемый плейлист, но и модель адаптации музыкального потока к текущей аудитории при сохранении приватности и контроля со стороны заведения. Тем самым обзор литературы задает общую исследовательскую траекторию: от широкого контекста рекомендаций и музыкальной среды к частной задаче групповых музыкальных рекомендаций и их прикладного сервисного воплощения.
